\section{Dualidad y kernels}

\begin{frame}{La formulación dual revela dónde entra el kernel}
\scriptsize
\begin{columns}[T,onlytextwidth]
\begin{column}{0.58\textwidth}
\begin{ccdebox}
\centering
\textbf{Problema dual del soft-margin SVM}
\[
\max_{\boldsymbol\alpha}
\quad
\sum_{i=1}^{n}\alpha_i
-\frac12\sum_{i=1}^{n}\sum_{j=1}^{n}
\alpha_i\alpha_j y_i y_j
\underbrace{\mathbf x_i^\top\mathbf x_j}_{\text{producto interno}}
\]
\[
\text{s.a.}\quad
0\le \alpha_i\le C,
\qquad
\sum_{i=1}^{n}\alpha_i y_i=0
\]
\end{ccdebox}
\vspace{0.12cm}
\begin{keybox}[fontupper=\footnotesize]
La optimización ya no necesita las coordenadas por separado: solo aparecen productos internos entre pares de observaciones.
\end{keybox}
\end{column}
\begin{column}{0.38\textwidth}
\begin{contrastbox}{ccdeaccent}{Consecuencia}
\footnotesize
Si reemplazamos
\[
\mathbf x_i^\top\mathbf x_j
\quad\text{por}\quad
K(\mathbf x_i,\mathbf x_j),
\]
obtenemos un clasificador lineal en otro espacio de características sin tener que construirlo explícitamente.
\end{contrastbox}
\end{column}
\end{columns}
\source{Hastie et al. (2009), cap. 12; Géron (2022), ``The Dual Problem'' y ``Kernelized SVMs''.}
\end{frame}

% ================================================================
\begin{frame}{Cambio de dimensionalidad: una frontera no lineal puede volverse lineal}
\small
\begin{columns}[c,onlytextwidth]
\begin{column}{0.38\textwidth}
\centering
\textbf{Espacio original: 1D}\\[0.15cm]
\begin{tikzpicture}[x=0.82cm,y=0.82cm]
  \draw[->,ccdemid] (-3.5,0)--(3.7,0) node[right]{\scriptsize $x$};
  \foreach \x in {-3,-2.4,2.4,3}{\fill[purplex] (\x,0) circle(3pt);}
  \foreach \x in {-0.9,-0.4,0,0.5,1.0}{\fill[ccdeaccent] (\x,0) circle(3pt);}
  \node[text=purplex,font=\scriptsize] at (0,1.0){clase $+1$ en ambos extremos};
  \node[text=ccdeaccent,font=\scriptsize] at (0,-0.9){clase $-1$ en el centro};
\end{tikzpicture}
\vspace{0.20cm}
\begin{contrastbox}{warn}{No separable con un umbral}
\scriptsize
Una sola frontera lineal en 1D es un punto; no puede dejar $+1$ en ambos extremos y $-1$ en el centro.
\end{contrastbox}
\end{column}
\begin{column}{0.12\textwidth}
\centering
\Huge\textcolor{ccdeaccent}{$\Longrightarrow$}\\[-0.05cm]
\scriptsize $\phi(x)=(x,x^2)$
\end{column}
\begin{column}{0.46\textwidth}
\centering
\textbf{Espacio transformado: 2D}\\[0.10cm]
\begin{tikzpicture}[x=0.70cm,y=0.60cm]
  \draw[->,ccdemid] (-3.6,0)--(3.7,0) node[right]{\scriptsize $x$};
  \draw[->,ccdemid] (0,-0.3)--(0,5.0) node[above]{\scriptsize $x^2$};
  \draw[ccdemid!60,thin,domain=-3.1:3.1,samples=80] plot (\x,{0.48*\x*\x});
  \foreach \x in {-3,-2.4,2.4,3}{\fill[purplex] (\x,{0.48*\x*\x}) circle(3pt);}
  \foreach \x in {-0.9,-0.4,0,0.5,1.0}{\fill[ccdeaccent] (\x,{0.48*\x*\x}) circle(3pt);}
  \draw[okgreen,very thick] (-3.3,2.35)--(3.3,2.35);
  \node[text=okgreen,font=\scriptsize\bfseries] at (2.0,2.70){hiperplano lineal};
\end{tikzpicture}
\vspace{0.08cm}
\begin{keybox}[fontupper=\scriptsize]
Una frontera lineal en el espacio ampliado se proyecta como una frontera \textbf{no lineal} en el espacio original.
\end{keybox}
\end{column}
\end{columns}
\source{Intuición de expansión de bases: James et al. (2013), sec. 9.3; Hastie et al. (2009), secs. 12.3 y 12.4.}
\end{frame}

% ================================================================
\begin{frame}{El kernel trick: calcular similitudes sin construir $\phi(x)$}
\footnotesize
\begin{columns}[T,onlytextwidth]
\begin{column}{0.52\textwidth}
\begin{ccdebox}
\centering
Supongamos una transformación potencialmente enorme:
\[
\phi:\mathbb R^p\to\mathcal H
\]
Un kernel válido satisface
\[
\boxed{K(\mathbf x,\mathbf z)=\langle \phi(\mathbf x),\phi(\mathbf z)\rangle_{\mathcal H}}
\]
La SVM solo necesita esos productos internos.
\end{ccdebox}
\vspace{0.14cm}
\begin{keybox}
No calculamos todas las nuevas variables. Construimos la matriz de Gram $K_{ij}=K(\mathbf x_i,\mathbf x_j)$ y optimizamos con ella.
\end{keybox}
\end{column}
\begin{column}{0.44\textwidth}
\begin{contrastbox}{ccdeaccent}{Por qué es potente}
\footnotesize
\begin{itemize}
  \item Un kernel polinómico equivale a usar interacciones y potencias.
  \item El RBF corresponde a un espacio de características implícito de dimensión infinita.
  \item El costo depende de evaluar $K$, no de materializar ese espacio.
\end{itemize}
\end{contrastbox}
\vspace{0.14cm}
\begin{contrastbox}{ccdeblue}{Condición matemática}
\footnotesize
Para kernels estándar, la matriz de Gram debe ser simétrica y semidefinida positiva; así puede interpretarse como un producto interno en algún espacio de Hilbert.
\end{contrastbox}
\end{column}
\end{columns}
\source{James et al. (2013), sec. 9.3; Hastie et al. (2009), secs. 12.3.3 y 12.3.7.}
\end{frame}

% ================================================================
\begin{frame}{Kernels más usados}
\scriptsize
\renewcommand{\arraystretch}{1.45}
\setlength{\tabcolsep}{5.0pt}
\begin{center}
\begin{tabular}{p{2.2cm}p{5.0cm}p{5.0cm}}
\toprule
\rowcolor{ccdenavy}
\textcolor{white}{\textbf{Kernel}} &
\textcolor{white}{\textbf{Fórmula}} &
\textcolor{white}{\textbf{Intuición / hiperparámetros}}\\
\midrule
Lineal & $K(x,z)=x^\top z$ & Sin expansión no lineal; suele funcionar bien cuando $p$ es grande.\\
\rowcolor{ccdeice!50}
Polinómico & $K(x,z)=(\gamma x^\top z+r)^d$ & $d$ controla grado; $\gamma$ escala la similitud; $r=\texttt{coef0}$ cambia peso entre términos de distinto orden.\\
RBF / Gaussiano & $K(x,z)=\exp(-\gamma\|x-z\|^2)$ & Similitud local. $\gamma$ alto: influencia muy local; $\gamma$ bajo: frontera más suave.\\
\rowcolor{ccdeice!50}
Sigmoide & $K(x,z)=\tanh(\gamma x^\top z+r)$ & Menos común; se parece formalmente a una neurona sigmoide.\\
\bottomrule
\end{tabular}
\end{center}
\vspace{0.16cm}
\begin{keybox}
En práctica, la comparación más importante suele ser \textbf{lineal vs. RBF}; los hiperparámetros deben elegirse con validación cruzada, no por ajuste visual al training set.
\end{keybox}
\source{James et al. (2013), Eqs. 9.22--9.24; Géron (2022), ``Polynomial Kernel'' y ``Gaussian RBF Kernel''.}
\end{frame}

% ================================================================
\begin{frame}{RBF en profundidad: $\gamma$ define el radio de influencia}
\small
\begin{columns}[T,onlytextwidth]
\begin{column}{0.66\textwidth}
\IfFileExists{fig_svm_gamma.png}{%
  \centering\includegraphics[width=\linewidth]{fig_svm_gamma.png}
}{%
  \begin{ccdebox}\centering\small Ejecuta \texttt{python codigo\_demo\_svm.py} para generar la figura de $\gamma$.\end{ccdebox}}
\end{column}
\begin{column}{0.30\textwidth}
\begin{contrastbox}{ccdeblue}{$\gamma$ bajo}
\scriptsize
La similitud cae lentamente con la distancia. Cada punto influye sobre una región grande: frontera suave, mayor sesgo.
\end{contrastbox}
\vspace{0.10cm}
\begin{contrastbox}{warn}{$\gamma$ alto}
\scriptsize
La similitud cae muy rápido. Cada punto influye localmente: frontera ondulada, mayor varianza.
\end{contrastbox}
\end{column}
\end{columns}
\source{Géron (2022), Figura 5-9; James et al. (2013), Eq. 9.24. Figura generada localmente.}
\end{frame}

% ================================================================
\begin{frame}{La función de decisión kernelizada}
\footnotesize
\begin{columns}[T,onlytextwidth]
\begin{column}{0.56\textwidth}
\begin{ccdebox}
\centering
Después de resolver el dual:
\[
f(\mathbf x)=
\sum_{i\in SV}\alpha_i y_i K(\mathbf x_i,\mathbf x)+b
\]
\[
\hat y=\operatorname{sign}(f(\mathbf x))
\]
\end{ccdebox}
\vspace{0.12cm}
\begin{keybox}
Solo aparecen los support vectors: aquellos con $\alpha_i>0$. Para un RBF, los support vectors cercanos a $x$ reciben mayor similitud y pesan más en la decisión.
\end{keybox}
\end{column}
\begin{column}{0.40\textwidth}
\begin{contrastbox}{ccdeaccent}{Lectura operacional}
\footnotesize
Para predecir una nueva observación:
\begin{enumerate}
  \item medir su similitud con cada support vector;
  \item ponderar por $\alpha_i y_i$;
  \item sumar el intercepto $b$;
  \item tomar el signo.
\end{enumerate}
\end{contrastbox}
\end{column}
\end{columns}
\source{James et al. (2013), Eq. 9.23; Hastie et al. (2009), cap. 12.}
\end{frame}

% ================================================================
